\section[Summary]{Summary}

\begin{frame}{Determinant Computation}
\begin{enumerate}
  \item Use $ad-bc$ for a $2\times2$ determinant.
  \item For cofactor expansion, choose a row or column containing many zeros.
  \item For a triangular matrix, multiply the diagonal entries.
  \item For a larger dense matrix, use elimination and track swaps and scalings.
\end{enumerate}
\end{frame}

\begin{frame}{Core Properties}
\[
\begin{array}{rcl}
R_i\leftarrow R_i+cR_j &:& \det\text{ unchanged},\\[.5ex]
R_i\leftrightarrow R_j &:& \det\text{ changes sign},\\[.5ex]
R_i\leftarrow kR_i &:& \det\text{ multiplied by }k,\\[.5ex]
\det(A^T)&=&\det A,\\[.5ex]
\det(AB)&=&(\det A)(\det B),\\[.5ex]
\det(cA)&=&c^n\det A.
\end{array}
\]
\end{frame}

\begin{frame}{Exit Check}
Let
\[
A=\begin{bmatrix}
2&1&0\\
4&3&1\\
0&5&2
\end{bmatrix}.
\]
\begin{enumerate}
  \item Which method would you choose to compute $\det A$?
  \item Which row operations would preserve the determinant?
  \item If two rows became equal, what could you conclude immediately?
\end{enumerate}
\pause
\begin{block}{Next lecture}
We will use determinants to study invertibility, linear systems, Cramer's rule, and geometric scaling.
\end{block}
\end{frame}

